\documentclass[tikz,border=8pt]{standalone}
\usepackage{tikz}
\usepackage{amsmath}
\usepackage{amssymb}
\usepackage{pgfplots}
\pgfplotsset{compat=1.18}
\usepackage{ctex}
\usetikzlibrary{calc, positioning, arrows.meta, decorations.pathreplacing}

\begin{document}
\begin{tikzpicture}

%% ===== 左图：原始单位网格 =====
\begin{axis}[
  name=axL,
  width=5.5cm, height=5.5cm,
  xmin=-1.5, xmax=2.5,
  ymin=-1.5, ymax=2.5,
  xlabel={$x$}, ylabel={$y$},
  title={原始 $\mathbb{R}^2$ 网格},
  title style={font=\small\bfseries},
  axis lines=center,
  xtick={-1,0,1,2},
  ytick={-1,0,1,2},
  xticklabel style={font=\scriptsize},
  yticklabel style={font=\scriptsize},
  grid=major, grid style={blue!30, dashed, thin},
  clip=false
]
  % 单位向量 e1=(1,0), e2=(0,1)
  \addplot[->,thick,blue!80!black,>=Stealth]
    coordinates {(0,0) (1,0)} node[below right,font=\small] {$\mathbf{e}_1$};
  \addplot[->,thick,blue!80!black,>=Stealth]
    coordinates {(0,0) (0,1)} node[above left,font=\small] {$\mathbf{e}_2$};
\end{axis}

%% 中间箭头
\node[font=\small] at ($(axL.east)+(0.7cm,0)$) {$\xrightarrow{\ A\ }$};

%% ===== 右图：变换后网格 =====
%% A = [[2,1],[1,2]], A*e1=(2,1), A*e2=(1,2)
%% 网格是整数组合 i*(2,1)+j*(1,2)
%% 范围 i,j in {-1,0,1,2}: 点从 (-3,-3) 到 (6,6)
\begin{axis}[
  name=axR,
  at={(axL.east)}, anchor=west, xshift=1.4cm,
  width=5.5cm, height=5.5cm,
  xmin=-0.5, xmax=5.5,
  ymin=-0.5, ymax=5.5,
  xlabel={$x$}, ylabel={$y$},
  title={变换后网格},
  title style={font=\small\bfseries},
  axis lines=center,
  xtick={0,1,2,3,4,5},
  ytick={0,1,2,3,4,5},
  xticklabel style={font=\scriptsize},
  yticklabel style={font=\scriptsize},
  grid=none,
  clip=false
]
  % 固定 i（沿 a2 方向的线），变化 j
  % i=-1: start=(-2+j*1, -1+j*2) j=-1..2 -> (-3,-3)...(0,3)
  \addplot[red!60, thin] coordinates {(-1,-1) (1,3)};
  % i=0: start=(j*1, j*2) j=-1..2 -> (-1,-2)...(2,4)
  \addplot[red!60, thin] coordinates {(0,0) (2,4)};
  % i=1: start=(2+j*1, 1+j*2) j=-1..2 -> (1,-1)..(4,5)
  \addplot[red!60, thin] coordinates {(1,0) (3,4)};  % clipped to visible
  % i=2: start=(4+j*1, 2+j*2) j=-1..2 -> (3,0)..(6,6)
  \addplot[red!60, thin] coordinates {(3,0) (5,4)};

  % 固定 j（沿 a1 方向的线），变化 i
  % j=-1: start=(i*2-1, i*1-2) i=-1..2 -> (-3,-3)..(3,0)
  \addplot[red!60, thin] coordinates {(0,-1) (4,1)};  % clipped
  % j=0: start=(i*2, i*1) i=-1..2 -> (-2,-1)..(4,2)
  \addplot[red!60, thin] coordinates {(0,0) (4,2)};
  % j=1: start=(i*2+1, i*1+2) i=-1..2 -> (-1,1)..(5,4)
  \addplot[red!60, thin] coordinates {(1,2) (5,4)};
  % j=2: start=(i*2+2, i*1+3) i=-1..2 -> (0,2)..(6,5)
  \addplot[red!60, thin] coordinates {(0,3) (4,5)};

  % 变换后基向量 A*e1=(2,1), A*e2=(1,2)
  \addplot[->,very thick,red!80!black,>=Stealth]
    coordinates {(0,0) (2,1)} node[below right,font=\small] {$A\mathbf{e}_1=(2,1)$};
  \addplot[->,very thick,red!80!black,>=Stealth]
    coordinates {(0,0) (1,2)} node[above right,font=\small] {$A\mathbf{e}_2=(1,2)$};
\end{axis}

%% 矩阵标注框（绝对位置，放在两图之间下方）
\node[font=\small, fill=white, draw=gray!50, thin, rounded corners=2pt,
      inner sep=5pt]
  at ($(axL.south east)!0.5!(axR.south west)+(0,-0.8cm)$)
  {$A=\begin{pmatrix}2&1\\1&2\end{pmatrix}$};

%% 整体标题
\node[font=\large\bfseries] at ($(axL.north)!0.5!(axR.north)+(0,1.2cm)$)
  {矩阵 $A$ 作为线性映射};

\end{tikzpicture}
\end{document}
